% TEMPLATE-MILC-v3.tex - plantilla maestra UNICA de las guias MILC v3.
%
% La usan las tres familias de guias, cada una con su driver:
%   · Grados 6-11  -> scripts/build-guias-g11.py
%   · Bebras       -> scripts/build-guias-bebras.py
%   · Semillero    -> scripts/build-guias-semillero.py
%
% Antes habia tres copias casi identicas (~400 lineas iguales, 6 distintas):
% cada mejora habia que aplicarla tres veces, y en la practica no se hacia
% (el motor de recursos vivia solo en dos de ellas). Lo que varia entre
% familias esta parametrizado en 6 placeholders que cada driver llena
% (se nombran aqui SIN su sintaxis real: el builder reemplaza por texto plano,
% tambien dentro de los comentarios, y un valor multilinea rompe el preambulo):
%
%   HEADER_IZQ        encabezado izquierdo de pagina
%   HEADER_DER        encabezado derecho de pagina
%   PORTADA_ETIQUETA  rotulo sobre el numero grande (GRADO/MOMENTO/MODULO)
%   PORTADA_SERIE     linea de serie de la portada
%   PORTADA_ID        identificador de la guia en la portada
%   PORTADA_CIRCULO   el circulito de grado (°), o vacio si no aplica
%
% El resto de claves las documenta content/guias/_SCHEMA.md. El script valida
% que no queden placeholders sin reemplazar antes de escribir el archivo final.

\documentclass[11pt,letterpaper]{article}
\usepackage[margin=2.0cm]{geometry}
\usepackage[table]{xcolor}
\usepackage{fontspec}
\usepackage[spanish]{babel}
\usepackage{tikz}
% Librerías TikZ para los diagramas de las guías (bloque `recursos:`):
%  · babel  → babel en español vuelve activos caracteres como «>», y TikZ los usa
%             en sus opciones (>=stealth, ->). Sin esta librería, cualquier
%             diagrama con flechas revienta con «\language@active@arg> has an
%             extra }». Debe ir ANTES de que se dibuje nada.
%  · positioning        → sintaxis «below left=2pt and 3pt».
%  · decorations.pathreplacing → llaves ({brace}) para señalar rangos.
\usetikzlibrary{babel,positioning,decorations.pathreplacing}
\usepackage{graphicx}
% Circuitos: el trabajo de electrónica se hace hoy con micro:bit y sus sensores
% integrados (MakeCode, sin cableado), así que no se necesitan esquemáticos.
% Si algún día entran montajes con protoboard, basta descomentar esta línea:
% los diagramas .tex/.tikz declarados en `recursos:` ya se incrustan solos.
% \usepackage{circuitikz}
\usepackage[most]{tcolorbox}
\usepackage{tabularx}
\usepackage{array}
\usepackage{fancyhdr}
\usepackage{titlesec}
\usepackage[document]{ragged2e}  % bandera a la derecha en todo el cuerpo
\usepackage{enumitem}
\usepackage{microtype}

% Helvetica Neue no trae la flecha U+2192, asi que cada «→» escrito literalmente
% en el YAML se compilaba a NADA, en silencio (462 flechas en 61 guias: rutas del
% tipo «paleta → tipografia → lockup» salian rotas). Mapeamos el caracter a su
% version matematica, que si tiene glifo. Asi el autor puede seguir escribiendo
% «→» comodamente en el YAML.
\usepackage{newunicodechar}
\newunicodechar{→}{\ensuremath{\rightarrow}}
% Mismo problema con los simbolos de marca y vinetas: el «✓» de «una palomita ✓»
% salia como cuadrito vacio en el PDF de todas las guias que lo usaban.
\usepackage{pifont}
\newunicodechar{✓}{\ding{51}}
\newunicodechar{✗}{\ding{55}}
\newunicodechar{✔}{\ding{52}}
\newunicodechar{•}{\textbullet}
\newunicodechar{≥}{\ensuremath{\geq}}
\newunicodechar{≤}{\ensuremath{\leq}}

\setmainfont{Helvetica Neue}
\setsansfont{Helvetica Neue}
\setlength{\parindent}{0pt}
\setlength{\parskip}{6pt}
% Sin particion de palabras: en 6.º-7.º una palabra cortada por guion al final
% de linea («Comuni-cacion») frena la lectura mucho mas que una linea corta.
\hyphenpenalty=10000
\exhyphenpenalty=10000
\pretolerance=2000
\tolerance=3000
\setlength{\emergencystretch}{3em}
\linespread{1.10}
\renewcommand{\arraystretch}{1.25}
\setlist{leftmargin=5mm,itemsep=1mm,topsep=1mm}
\newcolumntype{Y}{>{\RaggedRight\arraybackslash}X}

\definecolor{milcMagenta}{HTML}{E6007E}
\definecolor{milcVino}{HTML}{5A0038}
\definecolor{milcTurquesa}{HTML}{0093A5}
\definecolor{milcVerde}{HTML}{58A923}
\definecolor{milcMostaza}{HTML}{E5B400}
\definecolor{milcOcre}{HTML}{B86600}
\definecolor{milcNegro}{HTML}{241816}
\definecolor{milcGris}{HTML}{F7F7F4}
\definecolor{milcLinea}{HTML}{E8E2DD}
\definecolor{milcGradePrimary}{HTML}{FF2D87}
\definecolor{milcGradeDark}{HTML}{9F1B56}
\definecolor{milcGradeSoft}{HTML}{FFE0EE}
\definecolor{milcGradeAccent}{HTML}{CFE7FF}
\definecolor{milcGradeText}{HTML}{FFFFFF}
\color{milcNegro}

\pagestyle{fancy}
\fancyhf{}
\lhead{\small Tecnología e Informática · Grado 10}
\rhead{\small Guía 13 · MILC v3}
\cfoot{\small\thepage}
\renewcommand{\headrulewidth}{0pt}

\newtcolorbox{titlebox}[1]{
  enhanced,colback=#1,colframe=#1,arc=7mm,boxrule=0pt,
  left=7mm,right=7mm,top=3mm,bottom=3mm,
  before skip=8pt,after skip=8pt,
  fontupper=\sffamily\bfseries\Large\color{white}
}

\newtcolorbox{softbox}[3]{
  enhanced,colback=#2,colframe=#1,borderline west={4pt}{0pt}{#1},
  arc=2mm,boxrule=.4pt,left=4mm,right=4mm,top=3mm,bottom=3mm,
  before skip=6pt,after skip=8pt,title={#3},coltitle=white,
  colbacktitle=#1,fonttitle=\sffamily\bfseries,fontupper=\RaggedRight,
  attach boxed title to top left={xshift=3mm,yshift=-2mm},
  boxed title style={arc=2mm, boxrule=0pt}
}

\newtcolorbox{drawbox}[2]{
  enhanced,colback=white,colframe=#1,arc=4mm,boxrule=1pt,
  left=5mm,right=5mm,top=4mm,bottom=4mm,before skip=8pt,after skip=8pt,
  title={#2},coltitle=white,colbacktitle=#1,fonttitle=\sffamily\bfseries,
  fontupper=\RaggedRight,
  attach boxed title to top left={xshift=4mm,yshift=-2mm},
  boxed title style={arc=3mm, boxrule=0pt}
}

\newtcolorbox{infoband}[1]{
  enhanced,colback=#1!8,colframe=#1!50,arc=2mm,boxrule=.5pt,
  left=4mm,right=4mm,top=2mm,bottom=2mm,before skip=4pt,after skip=4pt,
  fontupper=\small\RaggedRight
}

% Puente narrativo entre secciones (contrato editorial Conéctate).
% Da continuidad: "Acabas de X. Ahora vas a Y porque Z."
% Visualmente: banda izquierda en color del grado, texto en itálica pequeña.
\newtcolorbox{puentebox}{
  enhanced,colback=milcGradeSoft,colframe=milcGradeDark,
  borderline west={3pt}{0pt}{milcGradeDark},
  arc=0mm,boxrule=0pt,left=5mm,right=4mm,top=2.5mm,bottom=2.5mm,
  before skip=4pt,after skip=8pt,
  fontupper=\itshape\small\color{milcNegro}\RaggedRight
}
% La flecha va en modo matematico a proposito: Helvetica Neue NO tiene el glifo
% U+2192, asi que \textrightarrow se compilaba a NADA («Missing character: There
% is no → in font Helvetica Neue») y el puente salia sin flecha. En modo
% matematico la flecha viene de la fuente matematica y si se dibuja.
\newcommand{\puente}[1]{%
  \begin{puentebox}%
    \textcolor{milcGradeDark}{$\rightarrow$}\hspace{2mm}#1%
  \end{puentebox}%
}

\newcommand{\linead}[1]{\noindent\color{milcLinea}\rule{#1}{0.5pt}\color{milcNegro}}
\newcommand{\respuesta}[1]{\par\vspace{2mm}\foreach \n in {1,...,#1}{\noindent\color{milcLinea}\rule{\linewidth}{0.5pt}\par\vspace{5mm}}\color{milcNegro}}
\newcommand{\checkbox}{\textcolor{milcGradeDark}{\fbox{\phantom{x}}}\hspace{5pt}}

\newcommand{\phasecard}[4]{
\begin{tcolorbox}[enhanced,colback=#3,colframe=#2,arc=3mm,boxrule=.5pt,height=3.05cm,valign=center,halign=center,left=1.5mm,right=1.5mm,top=2mm,bottom=2mm]
{\sffamily\bfseries\fontsize{22}{24}\selectfont\textcolor{#2}{#1}}\par
{\sffamily\bfseries\small #4}
\end{tcolorbox}
}

% ── Piezas del rediseno 2026 (legibilidad 6.º-11.º) ───────────────────
% Banda de estacion: reemplaza el titlebox macizo. Titulo a la izquierda,
% dato practico (tiempo/modalidad) a la derecha, en una sola linea.
\newcommand{\milcestacion}[2]{%
\begin{tcolorbox}[enhanced,colback=#1,colframe=#1,arc=2mm,boxrule=0pt,
  left=5mm,right=5mm,top=2.6mm,bottom=2.6mm,before skip=11pt,after skip=7pt]
{\sffamily\bfseries\fontsize{15}{18}\selectfont\color{white}#2}
\end{tcolorbox}}

% Tarjeta de paso numerada: sustituye las tablas «Pilar / Aplicacion», que
% eran formato de consulta docente, no de estudiante. El numero va en un
% circulo sobre el borde para que la secuencia se lea de un vistazo.
\newcommand{\milcpaso}[3]{%
\begin{tcolorbox}[enhanced,colback=white,colframe=milcLinea,arc=1.5mm,boxrule=.9pt,
  left=14mm,right=4mm,top=3mm,bottom=3mm,before skip=4pt,after skip=4pt,
  fontupper=\RaggedRight,
  overlay={\node[anchor=north west,circle,fill=#2,text=white,inner sep=0pt,
    minimum size=7.4mm,font=\sffamily\bfseries\small]
    at ([xshift=3.6mm,yshift=-3.2mm]frame.north west) {#1};}]
#3
\end{tcolorbox}}

% Casilla marcable de tamano util con lapiz (la anterior era \fbox{\phantom{x}},
% de alto variable segun el texto que la seguia).
\newcommand{\milcmarca}{\framebox[3.8mm]{\rule{0pt}{3.4mm}}}

% Fila marcable de lista suelta (checks de la estacion 1).
\newcommand{\milccheckrow}[1]{%
\par\vspace{1.8mm}\noindent
\makebox[6.4mm][l]{\raisebox{-.55mm}{\textcolor{milcNegro}{\milcmarca}}}%
\parbox[t]{\dimexpr\linewidth-6.4mm\relax}{\RaggedRight #1}\par}

% Celda marcable dentro de la rubrica (tabla de autoevaluacion). Misma
% sangria francesa, pero pensada para vivir dentro de una columna X.
\newcommand{\milcopcion}[1]{%
\makebox[5.2mm][l]{\raisebox{-.55mm}{\textcolor{milcNegro}{\milcmarca}}}%
\parbox[t]{\dimexpr\linewidth-5.2mm\relax}{\RaggedRight #1}}

% ── Recursos visuales de la guía ──────────────────────────────────────
% Declarados en el bloque `recursos:` del YAML y emitidos por el builder.
% \guiaFigura   → imagen rasterizada o PDF (foto, carta celeste, montaje).
% \guiaDiagrama → fuente TikZ/circuitikz (.tex) incrustada como vector:
%                 es la vía para circuitos y esquemáticos editables.
% #1 = ruta relativa al .tex · #2 = pie (puede ir vacío)
\newcommand{\guiaFigura}[2]{%
\begin{center}
\includegraphics[width=0.88\linewidth,height=0.38\textheight,keepaspectratio]{#1}\\[1.5mm]
{\footnotesize\itshape\textcolor{milcNegro!75}{#2}}
\end{center}
}

\newcommand{\guiaDiagrama}[2]{%
\begin{center}
\input{#1}\\[1.5mm]
{\footnotesize\itshape\textcolor{milcNegro!75}{#2}}
\end{center}
}

\begin{document}
\thispagestyle{empty}

% ── Portada ───────────────────────────────────────────────────────────
% Antes: un rectangulo de color a pagina completa. Se veia bien en pantalla,
% pero en la fotocopiadora del colegio se comia el toner y salia gris sucio.
% Ahora la banda ocupa el tercio superior y el resto va en blanco.
\begin{tikzpicture}[remember picture,overlay]
  \path[fill=milcGradePrimary]
    (current page.north west) rectangle ([yshift=-10.6cm]current page.north east);

  \node[anchor=north west,text=milcGradeText] at ([xshift=2.0cm,yshift=-1.55cm]current page.north west)
    {\sffamily\bfseries\fontsize{12}{14}\selectfont GRADO};
  \node[anchor=north west,text=milcGradeText,opacity=.85] at ([xshift=2.0cm,yshift=-2.35cm]current page.north west)
    {\sffamily\bfseries\fontsize{8.5}{10}\selectfont SERIE GUÍAS · TECNOLOGÍA E INFORMÁTICA};
  \node[anchor=north east,text=milcGradeText,opacity=.85,align=right] at ([xshift=-2.0cm,yshift=-1.55cm]current page.north east)
    {\sffamily\bfseries\fontsize{9}{12}\selectfont I.E. Sor María Juliana\\Cartago, Valle del Cauca};

  \node[anchor=west,text=milcGradeText] (gradonum) at ([xshift=1.85cm,yshift=-7.1cm]current page.north west)
    {\sffamily\bfseries\fontsize{112}{112}\selectfont 10};
  % El circulito de grado (°) solo aplica a las guias de grado («11°»); en
  % Bebras y Semillero el numero grande es un momento/modulo, no un grado.
  % Se posiciona relativo al nodo `gradonum`, no en coordenadas absolutas.
  \draw[milcGradeText,line width=1.6pt]
    ([xshift=2.0mm,yshift=-4.5mm]gradonum.north east) circle (.15cm);

  \node[anchor=west,text=milcGradeText,text width=12.3cm,align=left] at ([xshift=6.5cm,yshift=-6.6cm]current page.north west)
    {
     {\sffamily\bfseries\fontsize{9}{11}\selectfont\MakeUppercase{Período 2 · Informes técnicos y comunicación profesional con IA}}\\[3mm]
     {\sffamily\bfseries\fontsize{21}{25}\selectfont Google Docs ---\\estilos, índice\\y citas profesionales}\\[3.5mm]
     {\sffamily\bfseries\fontsize{15}{18}\selectfont Guía 13-10-TIC}
    };
\end{tikzpicture}

\vspace*{9.4cm}
\vfill

{\sffamily\bfseries\fontsize{8}{10}\selectfont\textcolor{milcGradeDark}{LO QUE TE LLEVAS HOY}}

\begin{tcolorbox}[enhanced,colback=white,colframe=milcNegro,arc=2mm,boxrule=1.6pt,
  left=5mm,right=5mm,top=4mm,bottom=4mm,before skip=3pt,after skip=12pt,
  fontupper=\RaggedRight\fontsize{11}{15.5}\selectfont]
Informe profesional en Google Docs (puede ser el informe técnico de la sesión 2 mejorado) aplicando los 5 elementos profesionales (a) estilos de título consistentes en 3 niveles, (b) tabla de contenidos autogenerada, (c) numeración automática de páginas, (d) al menos 3 citas con formato APA básico o IEEE, (e) al menos 1 imagen con pie de figura numerado + comparación lado a lado con un informe del mismo contenido sin estilos para demostrar la diferencia visual y funcional
\end{tcolorbox}

\begin{tcolorbox}[enhanced,colback=milcGris,colframe=milcGris,arc=2mm,boxrule=0pt,
  left=5mm,right=5mm,top=3.5mm,bottom=3.5mm,after skip=12pt]
{\sffamily\bfseries\fontsize{8}{10}\selectfont\textcolor{milcNegro!55}{ESTA GUÍA ES DE}}\\[3mm]
\begin{tabularx}{\linewidth}{X p{3.2cm} p{3.2cm}}
\linead{\linewidth} & \linead{3.0cm} & \linead{3.0cm}\\[-1mm]
{\fontsize{8.5}{10}\selectfont\textcolor{milcNegro!55}{Nombre completo}} &
{\fontsize{8.5}{10}\selectfont\textcolor{milcNegro!55}{Grupo}} &
{\fontsize{8.5}{10}\selectfont\textcolor{milcNegro!55}{Fecha}}\\
\end{tabularx}
\end{tcolorbox}

\vfill

\noindent\textcolor{milcLinea}{\rule{\linewidth}{1pt}}\\[2mm]
{\sffamily\bfseries\fontsize{9.5}{12}\selectfont Autor: Dr. Álvaro Cárdenas Orozco}\\[1mm]
{\sffamily\fontsize{8.5}{11}\selectfont\textcolor{milcNegro!55}{Metodología MILC · apertura ancestral · pensamiento computacional · triángulo Dussel–estoicismo–Floridi}}

\newpage

\section*{Apertura: Google Docs para informes profesionales --- estilos, índice y citas}

\begin{softbox}{milcGradeDark}{milcGradeSoft}{Saber ancestral}
Frente al juzgado del centro de Cartago, en la carrera 6 con calle 14, hubo durante décadas un personaje silencioso que sostenía la formalidad legal del pueblo: \textbf{el escribano público}. Cualquier vecino que necesitaba redactar carta, poder, petición o reclamo iba donde el escribano, que ponía mesa pequeña con máquina de escribir y carbón al frente. El escribano no inventaba el formato: usaba uno estricto y conocido. Cada tipo de documento tenía su \textbf{estructura inquebrantable}: las partes (quién y quién), los considerandos (los hechos), las cláusulas (los acuerdos), las firmas (la formalización). Si el documento se entregaba sin esa estructura, el juzgado lo rechazaba con la frase fría: \emph{``este documento no está bien hecho''}. La sabiduría era ancestral y precisa: \textbf{un documento legal mal organizado se cae solo}. El escribano memorizaba las estructuras de cada tipo: poder simple, poder amplio, petición, denuncia, carta de recomendación. Cada una con sus partes en orden estricto. Esa disciplina del oficio ancestral vive hoy en \textbf{los estilos y la estructura de Google Docs}: el documento profesional se sostiene en orden visible, no en improvisación.
\end{softbox}

\begin{softbox}{milcTurquesa}{milcGris}{Saber contemporáneo}
\textbf{Google Docs} es el procesador de texto gratuito más usado del mundo profesional. Su valor no está solo en escribir texto: está en sus \textbf{5 herramientas profesionales} que muchos usuarios casuales nunca aprovechan: (1) \textbf{Estilos de título}: \emph{Título principal}, \emph{Título 1}, \emph{Título 2}, \emph{Título 3}. Aplicarlos consistentemente da estructura visible y permite generar tabla de contenidos automática. (2) \textbf{Tabla de contenidos autogenerada}: menú Insertar → Tabla de contenidos. Se actualiza sola si cambias títulos o agregas secciones. (3) \textbf{Numeración automática de páginas}: menú Insertar → Número de página. Ubicación: arriba o abajo, centrado o lateral. (4) \textbf{Citas y bibliografía}: menú Herramientas → Citas. Permite formato APA, MLA, Chicago. Para grado 10, APA básico basta. (5) \textbf{Imágenes con pie de figura numerado}: insertar imagen → clic derecho → opciones de la imagen → texto alternativo + pie de figura debajo (\emph{Figura 1: descripción}). La regla profesional: \textbf{un documento que usa los 5 elementos se ve y se siente profesional; uno que no, parece tarea escolar}.
\end{softbox}

\begin{softbox}{milcMostaza}{milcGris}{Pregunta-puente}
\textit{¿Qué sabía el escribano del juzgado al respetar la estructura inquebrantable de cada documento legal, que el novato olvida cuando escribe informe en Docs sin aplicar los estilos? ¿Y por qué los 5 elementos profesionales de Docs hacen tanta diferencia visual y funcional?}
\end{softbox}

\begin{softbox}{milcVerde}{milcGris}{Saber y saber hacer}
Al terminar podrás: (1) \textbf{identificar} los 5 elementos profesionales de Google Docs y cuándo se aplican; (2) \textbf{aplicar} estilos de título consistentes, tabla de contenidos, numeración, citas APA básicas, imagen con pie de figura; (3) \textbf{analizar} las diferencias visuales y funcionales entre un documento con estilos y uno sin; (4) \textbf{evaluar} tu informe técnico anterior aplicándole los 5 elementos para elevarlo a nivel profesional.
\end{softbox}



\small
\begin{tabularx}{\textwidth}{>{\bfseries}p{3.0cm}Y}
\rowcolor{milcGradePrimary!12}Autor & Dr. Álvaro Cárdenas Orozco\\
DBA & Comunicación profesional con asistencia de IA y herramientas ofimáticas gratuitas (MEN, T\&I)\\
Referentes & Markdown como estándar abierto · Google Workspace · Estoicismo (claridad)\\
Producto final & Informe profesional en Google Docs (puede ser el informe técnico de la sesión 2 mejorado) aplicando los 5 elementos profesionales (a) estilos de título consistentes en 3 niveles, (b) tabla de contenidos autogenerada, (c) numeración automática de páginas, (d) al menos 3 citas con formato APA básico o IEEE, (e) al menos 1 imagen con pie de figura numerado + comparación lado a lado con un informe del mismo contenido sin estilos para demostrar la diferencia visual y funcional\\
\end{tabularx}
\normalsize

\puente{Antes de empezar, mira el viaje completo. En la sesión 2 escribiste informe técnico con estructura. Hoy le pones \textbf{forma profesional} en Google Docs. Las próximas sesiones (Markdown, encuestas, análisis con IA, informe comercial, correo profesional, mini-proyecto, sustentación) usarán Docs como herramienta principal.}

\milcestacion{milcGradeDark}{Ruta MILC: así vamos a aprender}
\begin{tabularx}{\textwidth}{>{\centering\arraybackslash}X>{\centering\arraybackslash}X>{\centering\arraybackslash}X>{\centering\arraybackslash}X}
\phasecard{1}{milcVerde}{milcGris}{Escucha\\[-1mm]\small Observo, escucho y pregunto.} &
\phasecard{2}{milcTurquesa}{milcGris}{Entender\\[-1mm]\small Sistematizo y ordeno.} &
\phasecard{3}{milcMagenta}{milcGris}{Hacer\\[-1mm]\small Construyo y aplico.} &
\phasecard{4}{milcVino}{milcGris}{Cierre\\[-1mm]\small Evalúo y reflexiono.}
\end{tabularx}


\puente{Antes de teorizar, vas a hacer un \emph{experimento rápido}: abre Google Docs, escribe 5 líneas con \emph{``Hola, esto es un título''} y debajo \emph{``Este es el cuerpo''}. Pruébalo sin aplicar estilos. Después aplica \emph{Título 1} a la primera línea y \emph{Cuerpo} a la segunda. Compara las 2 versiones visualmente.}

\milcestacion{milcVerde}{Estación 1 de 3 · Escucha}

\begin{softbox}{milcVerde}{milcGris}{Escena para escuchar}
\textbf{Actividad 1 · IDENTIFICA --- Experimento del estilo} (15 min · individual con dispositivo). Abre Google Docs en un documento nuevo. Escribe 5 líneas: (1) un título principal (\emph{``Informe sobre el recreo''}), (2) un subtítulo (\emph{``Análisis del tiempo disponible''}), (3) una línea de cuerpo, (4) otro subtítulo (\emph{``Hallazgos''}), (5) más cuerpo. Primero deja todo el texto en estilo \emph{Texto normal}. Toma captura. Después aplica: línea 1 → \emph{Título principal}, líneas 2 y 4 → \emph{Título 1}, líneas 3 y 5 → \emph{Texto normal}. Toma segunda captura. Compara las 2 visualmente.
\end{softbox}

{\sffamily\bfseries\fontsize{8}{10}\selectfont\textcolor{milcNegro!55}{MARCA CUANDO LO HAYAS HECHO}}
\milccheckrow{Hice el experimento de los 5 líneas en Google Docs.}
\milccheckrow{Tomé captura antes y después de aplicar estilos.}
\milccheckrow{Comparé visualmente las 2 versiones.}

\begin{infoband}{milcVerde}


\textbf{Lo que va al cuaderno (Actividad 1 · IDENTIFICA).} Título: «Actividad 1 --- Experimento del estilo». \textbf{Formato:} 2 capturas + observación de la diferencia. \textbf{Sabes que terminaste cuando} ves que los estilos no son solo cosmética, son estructura.
\end{infoband}

\puente{Ya viste el efecto. Ahora vas a entender la \textbf{anatomía profesional} de los 5 elementos de Google Docs y los 5 errores típicos del usuario casual.}

\milcestacion{milcTurquesa}{Estación 2 de 3 · Entender}

\begin{softbox}{milcTurquesa}{milcGris}{Lo que vas a entender}
Tu experimento te mostró que los estilos cambian la apariencia. Ahora necesitas la \textbf{anatomía profesional} de los 5 elementos: cómo aplicarlos, qué resuelve cada uno, los errores típicos.
\end{softbox}

\milcpaso{1}{milcTurquesa}{Los \textbf{5 elementos profesionales} de Google Docs se descomponen así: (1) \textbf{Estilos de título}: menú desplegable en la barra de herramientas (donde dice \emph{Texto normal}). Elige Título 1 para capítulos principales, Título 2 para subsecciones, Título 3 para sub-subsecciones. Aplica consistentemente. (2) \textbf{Tabla de contenidos}: menú Insertar → Tabla de contenidos. Aparece automáticamente con los títulos aplicados. Botón \emph{actualizar} en la esquina si cambias estructura. (3) \textbf{Numeración de páginas}: menú Insertar → Números de página → elegir posición. (4) \textbf{Citas APA básicas}: menú Herramientas → Citas → APA. Permite agregar referencias de libros, artículos, sitios web. Genera bibliografía al final. (5) \textbf{Imagen con pie de figura}: insertar imagen → debajo, escribir \emph{Figura 1: descripción} en estilo \emph{Subtítulo} o cuerpo en cursiva. Numerar las figuras secuencialmente.}
\milcpaso{2}{milcTurquesa}{El \textbf{formato APA básico} para citas tiene 3 elementos mínimos: (a) Apellido del autor + año entre paréntesis dentro del texto (\emph{``Como mostró González (2024)...''}). (b) Lista de referencias al final del documento con: Apellido, Nombre. (año). Título. Editorial o URL. (c) Para sitios web: Nombre del sitio o autor. (año). Título de la página. URL. Recuperado el día/mes/año. La regla profesional: \textbf{toda afirmación que no sea de tu cabeza debe tener cita}.}
\milcpaso{3}{milcTurquesa}{Lo esencial cabe en una pregunta: \emph{¿este documento sería utilizable por alguien que solo tiene 2 minutos para mirar?}. Si el lector apurado puede mirar tabla de contenidos, saltar a la sección que le interesa, y leer el resumen de esa sección, el documento es profesional. Si tiene que leer todo de corrido sin orientación visual, falla en estructura.}
\milcpaso{4}{milcTurquesa}{Algoritmo: (1) abrir tu informe técnico de la sesión 2 en Docs. (2) aplicar \emph{Título principal} al título del informe. (3) aplicar \emph{Título 1} a Introducción, Desarrollo, Conclusiones. (4) aplicar \emph{Título 2} a subsecciones (Contexto, Objetivo, Alcance, etc.). (5) generar tabla de contenidos. (6) agregar numeración de páginas. (7) buscar 3 fuentes que sostengan afirmaciones del informe y agregar citas APA. (8) insertar al menos 1 imagen relevante con pie de figura.}

\begin{drawbox}{milcTurquesa}{Anatomía de Google Docs profesional}
\textbf{Estilos:} Título principal + Título 1, 2, 3. \\[1mm] \textbf{TOC:} Insertar → Tabla de contenidos. \\[1mm] \textbf{Numeración:} Insertar → Número de página. \\[1mm] \textbf{Citas APA:} Herramientas → Citas → APA. \\[1mm] \textbf{Imagen + pie:} Insertar imagen + \emph{Figura 1: descripción}.
\end{drawbox}

\begin{softbox}{milcMostaza}{milcGris}{Errores comunes y su corrección}
(1) \textbf{Formatear a mano sin estilos}: usar \emph{negrita + tamaño grande} en vez de \emph{Título 1}. (2) \textbf{Inconsistencia de estilos}: algunos títulos en negrita manual, otros con estilo. (3) \textbf{Sin tabla de contenidos}: el lector no puede navegar. (4) \textbf{Citas pegadas a mano}: en lugar de usar herramientas → Citas. (5) \textbf{Imágenes sin pie de figura}: no se sabe qué representan ni se pueden referenciar en el texto.
\end{softbox}

\begin{infoband}{milcTurquesa}


\textbf{Lo que va al cuaderno (Actividad 2 · APLICA).} Título: «Actividad 2 --- Anatomía Docs profesional». \textbf{Formato:} (a) los 5 elementos con su menú; (b) plantilla básica APA; (c) los 5 errores con marca del que más temes. \textbf{Sabes que terminaste cuando} sabes dónde está cada herramienta sin abrir esta guía.
\end{infoband}

\puente{Tienes el método. Toca aplicarlo: tomas el informe técnico de la sesión 2 y lo elevas a versión profesional aplicando los 5 elementos.}

\milcestacion{milcMagenta}{Estación 3 de 3 · Hacer}

\begin{softbox}{milcMagenta}{milcGris}{Cómo vas a construir tu producto}
\textbf{Actividad 3 · ANALIZA + EVALÚA --- Informe profesional + comparación} (40 min · individual con dispositivo). Hora de aplicar. Tomas el informe técnico de la sesión 2 y le aplicas los 5 elementos profesionales. Comparas con la versión sin estilos.
\end{softbox}

\milcpaso{1}{milcMagenta}{Tu informe profesional se construye en 6 pasos: (1) duplicar el informe técnico de la sesión 2 en Docs (uno será la versión sin estilos, otro será la versión profesional). (2) en la versión profesional, aplicar \emph{Título principal} al título del documento. (3) aplicar \emph{Título 1} a las 3 secciones principales y \emph{Título 2} a subsecciones. (4) generar tabla de contenidos al inicio. (5) agregar numeración de páginas, 3 citas APA, 1 imagen con pie de figura. (6) hacer comparación lado a lado con la versión sin estilos.}
\milcpaso{2}{milcMagenta}{Sigue el patrón \textbf{``estilos consistentes en todo el documento''}: si un capítulo usa Título 1 pero el siguiente usa negrita manual, el documento se rompe. La phronesis del editor consiste en \textbf{aplicar la misma regla a todos los elementos del mismo nivel}.}
\milcpaso{3}{milcMagenta}{La comparación lado a lado se hace tomando capturas de pantalla de ambas versiones del mismo informe. En 1 página coloca: versión sin estilos a la izquierda, versión profesional a la derecha. Anota 3 diferencias visuales evidentes. Esa comparación demuestra que los estilos no son cosmética; son comunicación.}
\milcpaso{4}{milcMagenta}{Algoritmo: (1) duplicar informe. (2) aplicar estilos. (3) TOC. (4) numeración. (5) 3 citas APA. (6) imagen + pie. (7) comparación con versión sin estilos. (8) verificar que la TOC se actualizó correctamente y la numeración funciona.}

\begin{drawbox}{milcMagenta}{Checklist antes de entregar}
\checkbox Estilos de título aplicados consistentemente. \\[1mm] \checkbox Tabla de contenidos autogenerada al inicio. \\[1mm] \checkbox Numeración de páginas automática. \\[1mm] \checkbox 3+ citas APA con referencias al final. \\[1mm] \checkbox 1+ imagen con pie de figura numerado. \\[1mm] \checkbox Comparación lado a lado con versión sin estilos. \\[1mm] \checkbox 3 diferencias visuales identificadas.
\end{drawbox}

\begin{softbox}{milcMagenta}{milcGris}{Plantilla de trabajo}
\textbf{Estilos.} \textit{Título principal + Título 1, 2, 3 aplicados.} \\[1mm] \textbf{TOC.} \textit{Autogenerada con menú.} \\[1mm] \textbf{Numeración.} \textit{Insertar Número de página.} \\[1mm] \textbf{Citas APA.} \textit{3 referencias con Herramientas Citas.} \\[1mm] \textbf{Imagen.} \textit{Con pie Figura 1 descripción.} \\[1mm] \textbf{Comparación.} \textit{Sin estilos vs profesional.}
\end{softbox}

\begin{infoband}{milcMagenta}


\textbf{Lo que va al cuaderno (Actividad 3 · ANALIZA + EVALÚA).} Título: «Actividad 3 --- Informe profesional en Docs». \textbf{Formato:} (a) link al documento; (b) capturas comparativas; (c) 3 diferencias identificadas. \textbf{Sabes que terminaste cuando} podrías mostrar el documento a un docente y se vería como pieza profesional.
\end{infoband}

\puente{Lo que sigue resume \emph{qué} entregas hoy: informe profesional con los 5 elementos + comparación lado a lado con versión sin estilos. El criterio principal: que la diferencia visual y funcional sea evidente al primer vistazo.}

\milcestacion{milcMostaza}{Producto de la sesión}
\textbf{Informe profesional en Docs + comparación lado a lado + análisis}.
\begin{softbox}{milcMostaza}{milcGris}{Criterios de calidad}
(1) 5 elementos profesionales aplicados. (2) Estilos consistentes en todo el documento. (3) TOC autogenerada y funcional. (4) Citas APA con referencias al final. (5) Imagen con pie de figura. (6) Comparación clara con versión sin estilos.
\end{softbox}

\puente{Antes de cerrar, mira los estilos desde las cinco dimensiones humanas. Aplicar los 5 elementos es respeto por el lector profesional; entregar texto crudo es desperdicio del oficio.}

\milcestacion{milcVino}{Evaluación liberadora · cinco dimensiones}

\begin{softbox}{milcVino}{milcGris}{Sentido de la evaluación}
Evaluar no es solo revisar una nota. Es reconocer qué aprendí, qué cambió en mí, qué emociones manejé, cómo me relaciono con la ciudad y la comunidad, y qué le dejaré a quien viene detrás.
\end{softbox}

\textbf{Mi reflexión liberadora en cinco dimensiones.} Completa cada frase iniciada:

\small
\begin{tabularx}{\textwidth}{>{\bfseries\RaggedRight\arraybackslash}p{3.4cm}Y>{\RaggedRight\arraybackslash}p{4.6cm}}
\rowcolor{milcVino}\color{white}Dimensión & \color{white}\bfseries Mi respuesta (frame starter) & \color{white}\bfseries Algo que descubrí\\
1. Personal & Lo que más cambió en mí fue \linead{6.5cm} & \linead{4.2cm}\\
2. Emocional & La emoción más fuerte fue \linead{2.5cm} y la regulé \linead{3cm} & \linead{4.2cm}\\
3. Ciudadana & En democracia esto importa porque \linead{6cm} & \linead{4.2cm}\\
4. Local (Cartago) & En mi barrio sería útil para \linead{6.5cm} & \linead{4.2cm}\\
5. Intergeneracional & A mi abuela/abuelo le diría \linead{6.5cm} & \linead{4.2cm}\\
\end{tabularx}
\normalsize

\begin{infoband}{milcVino}
\textbf{Para la próxima sustentación.} Vuelve a esta tabla en seis meses. Verás que las respuestas cambian — no porque lo aprendido cambie, sino porque tú cambias. La evaluación liberadora es una conversación contigo a través del tiempo.
\end{infoband}


\puente{Y para terminar, tres voces sobre el orden documental. Dussel sobre los documentos que respetan al lector profesional, Marco Aurelio sobre la disciplina de los estilos, Floridi sobre la documentación bien estructurada como ética del oficio.}

\milcestacion{milcGradeDark}{Triángulo de pensamiento · cierre filosófico}

\begin{softbox}{milcVino}{milcGris}{Dussel · lente del nosotros}
\textit{Un documento estructurado respeta al lector profesional; uno desordenado le quita tiempo y dignidad.}

Quien recibe tu informe profesional (docente, coordinador, futuro jefe) tiene poco tiempo. La filosofía de la liberación pide ese cuidado: \textbf{estructurar visualmente para que el lector apurado pueda navegar}. La TOC permite saltar a la sección clave; los estilos permiten escanear el documento; las citas permiten verificar. La phronesis del oficio consiste en \textbf{servir al lector con orden visible}.

\textbf{¿Mi documento sirve al lector profesional con orden visible, o le exige leer todo de corrido?}
\end{softbox}
\linead{\linewidth}\\[1mm]
\linead{\linewidth}

\begin{softbox}{milcMostaza}{milcGris}{Estoicismo · Marco Aurelio · cuidado interior}
\textit{Aplicar estilos con disciplina es virtud profesional; formatear a mano por capricho es vanidad.}

Es tentador formatear a mano (negrita aquí, tamaño 14 allá) porque parece más rápido. La sabiduría estoica recomienda lo opuesto: \emph{usa los estilos del sistema con disciplina}. La uniformidad que dan los estilos es virtud profesional. Quien formatea a mano produce documentos que se rompen al menor cambio.

\textbf{¿Estoy usando los estilos del sistema, o formateando a mano por costumbre?}
\end{softbox}
\linead{\linewidth}\\[1mm]
\linead{\linewidth}

\begin{softbox}{milcTurquesa}{milcGris}{Floridi · lente de la infoesfera}
\textit{La documentación bien estructurada es ética profesional en la era de los documentos digitales compartidos.}

En la era digital, los documentos circulan: se comparten con compañeros, se exportan a PDF, se editan en grupo. La ética profesional pide que esa circulación funcione: estilos que se mantengan, TOC que se actualice, citas que se puedan verificar. Tu documento de hoy te entrena en esa práctica que rige toda profesión que produzca documentos.

\textbf{¿Mi documento funciona bien al circular, o se rompe en cuanto alguien más lo abre?}
\end{softbox}
\linead{\linewidth}\\[1mm]
\linead{\linewidth}

\begin{infoband}{milcNegro}
\textbf{Nota para el docente.} Las tres formulaciones del triángulo son la síntesis del autor de la guía, no citas textuales de los pensadores. Si vas a citarlos en clase, remite a la obra original.
\end{infoband}

\milcestacion{milcVino}{Mi autoevaluación · marca una casilla por fila}

{\small
\setlength{\extrarowheight}{2.2pt}
\begin{tabularx}{\textwidth}{>{\RaggedRight\arraybackslash\bfseries}p{3.0cm}YYY}
\rowcolor{milcVino}\color{white}\bfseries Criterio & \color{white}\bfseries Lo logré & \color{white}\bfseries En proceso & \color{white}\bfseries Necesito apoyo\\[.6mm]
Comprensión del tema & \milcopcion{Explico las ideas con mis palabras.} & \milcopcion{Reconozco las ideas, falta precisión.} & \milcopcion{Necesito repasar lo básico.}\\[.8mm]
\rowcolor{milcGris}Pensamiento computacional & \milcopcion{Usé los pilares al armar mi producto.} & \milcopcion{Usé algunos pilares.} & \milcopcion{Necesito releer la estación 2.}\\[.8mm]
Aplicación y producto & \milcopcion{Mi producto cumple los criterios.} & \milcopcion{Está armado, con ajustes pendientes.} & \milcopcion{Aún está en construcción.}\\[.8mm]
\rowcolor{milcGris}Reflexión 5 dimensiones & \milcopcion{Respondí cada una con honestidad.} & \milcopcion{Respondí algunas con detalle.} & \milcopcion{Mis respuestas son muy generales.}\\[.8mm]
Triángulo de pensamiento & \milcopcion{Conecté las 3 voces con mi trabajo.} & \milcopcion{Conecté al menos una.} & \milcopcion{No las relaciono con mi proceso.}\\[.8mm]
\end{tabularx}}

\vspace{3mm}
\textbf{Hoy entendí que:}\\[1mm]
\linead{\linewidth}\\[1mm]
\linead{\linewidth}

\textbf{Mi compromiso para la próxima sesión es:}\\[1mm]
\linead{\linewidth}\\[1mm]
\linead{\linewidth}

\begin{softbox}{milcMostaza}{milcGris}{Cita de despedida · Marco Aurelio}
\textit{``El obstáculo en el camino se vuelve el camino. Lo que estorba al camino se vuelve camino.''}\\[1mm]
Cada error de hoy, bien observado, se convierte en pieza del camino que pisarás mañana.
\end{softbox}

\small
\begin{tabularx}{\textwidth}{>{\bfseries}p{3.6cm}Y}
\rowcolor{milcGradeDark!12}Nombre completo: & \linead{10cm}\\
Fecha: & \linead{4cm} \quad Curso/grupo: \linead{4cm}\\
Mi firma: & \linead{9cm}\\
\end{tabularx}
\normalsize

\begin{infoband}{milcGradeDark}
\textbf{Para la próxima sesión.} Llega con tu informe profesional en Docs y la comparación. En la sesión 4 vas a aprender \textbf{Markdown}: el formato portable que se puede convertir a PDF, HTML, libro, sin perder estructura. Docs es flexible; Markdown es portable. Cada uno tiene su lugar.
\end{infoband}

\end{document}
